\documentclass[12pt,a4paper]{article}

\usepackage{fontawesome5}
\providecommand{\faPencilAlt}{\faIcon{pencil-alt}}

\usepackage{styles/couleurs}
\usepackage{styles/mise_en_page}
\usepackage{styles/boites}
\usepackage{styles/maths}
\usepackage{styles/figures}

\usepackage{enumitem}
\usepackage{array}
\usepackage{tabularx}
\usepackage{amsmath}
\usepackage{amssymb}
\usepackage{tikz}
\usepackage{xparse}

%==================================================
% Réglages
%==================================================

\setlist[itemize]{
  leftmargin=1.4em,
  itemsep=0.08em,
  topsep=0.12em
}

\setlist[enumerate]{
  leftmargin=*,
  itemsep=0.12em,
  topsep=0.12em
}

\newcommand{\titreseq}[1]{%
\begin{center}
  {\sffamily\LARGE\bfseries\color{TitreC}
  Séquence 3 -- Nombres relatifs}\\[0.1em]
  {\sffamily\large\bfseries\color{SavoirC}
  #1}
\end{center}
}

\begin{document}

\titreseq{Produits, quotients et enchaînements}

%==================================================
\section*{1. Rappel : distributivité et nature d'une expression}
%==================================================

Pour reconnaître la nature d'une expression, on regarde
\textbf{l'opération effectuée en dernier}, en tenant compte
des parenthèses et des priorités opératoires.

\begin{center}
\renewcommand{\arraystretch}{1.25}

\begin{tabularx}{0.78\textwidth}{
|>{\centering\arraybackslash}X
|>{\centering\arraybackslash}X|}
\hline

\textbf{Dernière opération} & \textbf{Nature}\\
\hline
addition & somme\\
\hline
soustraction & différence\\
\hline
multiplication & produit\\
\hline
division & quotient\\
\hline

\end{tabularx}
\end{center}

\aretenir{
\begin{itemize}
  \item Une expression \textbf{développée} est écrite sous la forme
  d'une somme ou d'une différence.
  \item Une expression \textbf{factorisée} est écrite sous la forme
  d'un produit.
\end{itemize}
}

\exemple{
\[
7\times14+7\times6
=
7\times(14+6)
\]

À gauche, l'expression est développée ; à droite, elle est factorisée.
}

\aretenir{
\begin{itemize}
  \item \textbf{Développer} : transformer un produit
  en une somme ou une différence.
  \item \textbf{Factoriser} : transformer une somme ou une différence
  en un produit.
\end{itemize}
}

%==================================================
\section*{2. Multiplier deux nombres relatifs}
%==================================================

\methode{
Pour multiplier deux nombres relatifs :

\begin{enumerate}
  \item on détermine le \textbf{signe} du résultat ;
  \item on multiplie les \textbf{distances à zéro}.
\end{enumerate}
}

\begin{center}
\renewcommand{\arraystretch}{1.3}

\begin{tabularx}{0.78\textwidth}{
|>{\centering\arraybackslash}X
|>{\centering\arraybackslash}X|}
\hline

\textbf{Signes des facteurs} & \textbf{Signe du produit}\\
\hline
même signe & positif\\
\hline
signes différents & négatif\\
\hline

\end{tabularx}
\end{center}

\exemple{
\[
(-3)\times(+7)=-21
\qquad
(-2{,}5)\times(-4)=10
\]
}

%==================================================
\section*{3. Produit de plusieurs facteurs}
%==================================================

\propriete{
Pour un produit de plusieurs facteurs non nuls :

\begin{itemize}
  \item un nombre \textbf{pair} de facteurs négatifs
  donne un produit positif ;
  \item un nombre \textbf{impair} de facteurs négatifs
  donne un produit négatif.
\end{itemize}

Si un facteur est nul, le produit est nul.
}

\exemple{
\[
(-2)\times3\times(-5)=30
\qquad
(-2)\times(-3)\times(-4)=-24
\]
}

%==================================================
\section*{4. Diviser deux nombres relatifs}
%==================================================

\methode{
Pour diviser deux nombres relatifs, on utilise la
\textbf{même règle de signes} que pour la multiplication :

\begin{itemize}
  \item même signe $\rightarrow$ quotient positif ;
  \item signes différents $\rightarrow$ quotient négatif.
\end{itemize}

Puis on divise les distances à zéro.
}

\exemple{
\[
(-42)\div(+6)=-7
\qquad
(-35)\div(-5)=7
\]
}

\aretenir{
On ne peut pas diviser par zéro.
}

%==================================================
\section*{5. Enchaînements de calculs}
%==================================================

\propriete{
Dans une expression :

\begin{enumerate}
  \item on commence par les calculs entre parenthèses ;
  \item on effectue les multiplications et les divisions ;
  \item on termine par les additions et les soustractions.
\end{enumerate}

Lorsque plusieurs opérations ont la même priorité,
on les effectue de gauche à droite.
}

%--------------------------------------------------
% Comparaison bonne / mauvaise rédaction
%--------------------------------------------------

\begin{tcolorbox}[
  enhanced,
  colback=white,
  colframe=NeutreC!45,
  boxrule=0.6pt,
  arc=2mm,
  left=3mm,
  right=3mm,
  top=2mm,
  bottom=2mm,
  before skip=3pt,
  after skip=3pt
]

\noindent
\begin{minipage}[t]{0.475\linewidth}

{\sffamily\bfseries\color{SavoirC}
Une rédaction correcte}

\smallskip

On effectue les opérations prioritaires et on recopie le reste :

\[
\begin{aligned}
A&=(-24)\div(-6)+5\times(-2)\\
 &=4+(-10)\\
 &=-6
\end{aligned}
\]

Chaque ligne a la \textbf{même valeur} que la précédente.

\end{minipage}
\hfill
\begin{minipage}[t]{0.475\linewidth}

{\sffamily\bfseries\color{ActionC}
Ce qu'il ne faut pas faire}

\smallskip

\[
\begin{aligned}
A&=(-24)\div(-6)+5\times(-2)\\
 &=4\\
 &=5\times(-2)\\
 &=-10\\
 &=4+(-10)\\
 &=-6
\end{aligned}
\]

Ici, le signe \(=\) est utilisé comme s'il signifiait « puis ».

Or :
\[
4\neq5\times(-2)
\]

Le signe \(=\) relie deux expressions de
\textbf{même valeur}.

\end{minipage}

\end{tcolorbox}

%--------------------------------------------------
% Synthèse finale compacte
%--------------------------------------------------

\aretenir{
À chaque nouvelle ligne, on effectue les opérations prioritaires,
on recopie le reste de l'expression et on conserve une
\textbf{égalité vraie}.

\medskip

Avant ou après un calcul, on peut vérifier le \textbf{signe},
l'\textbf{ordre de grandeur} et la \textbf{cohérence du résultat}.

\medskip

Les parenthèses sont surtout utiles pour éviter un doute
entre deux signes qui se suivent :
\[
5\times(-3)
\qquad
12\div(-4)
\]
}

\end{document}